(1) 支点匀加速系附加惯性力为 \(-ma\hat x\)。令 \(x=\ell\sin\theta\)、\(y=-\ell\cos\theta\)，有效势差可写为
\[
U_{\rm eff}(\theta)=m\ell[g(1-\cos\theta)+a\sin\theta].
\]
动力学方程 \(m\ell^2\ddot\theta=-\partial_\theta U_{\rm eff}\)，即
\[
\ddot\theta+\frac g\ell\sin\theta+\frac a\ell\cos\theta=0.
\]
在匀加速系内 \(E_{\rm eff}=m\ell^2\dot\theta^2/2+U_{\rm eff}\) 守恒，虽然实验室中的机械能不必守恒。

(2) 令 \(g_*=\sqrt{g^2+a^2}\)，\(\varphi=\arctan(a/g)\in(0,\pi/2)\)，则 \(g\sin\theta+a\cos\theta=g_*\sin(\theta+\varphi)\)。平衡角为 \(\theta_s=-\varphi\pmod{2\pi}\) 与 \(\theta_u=\pi-\varphi\pmod{2\pi}\)。有效势曲率分别为 \(+m\ell g_*\) 与 \(-m\ell g_*\)，故前者稳定、后者不稳定。近稳定点频率 \(\omega_s=\sqrt{g_*/\ell}\)。

(3) 突然开始加速度后，角度和角速度不能瞬时跳变，初值为 \(\theta(0)=0,\dot\theta(0)=0\)。有效势等价于 \(mg_*\ell[1-\cos(\theta+\varphi)]\) 加常数。其初始相对稳定位置的偏角为 \(\varphi\)，故摆动的精确角振幅为 \(\varphi\)（但有限振幅运动不是简谐）。相邻转向点在 \(\theta=0\) 与 \(\theta=-2\varphi\)，无阻力时往返等高。

(4) 对质点在非惯性系的向心径向方程投影，杆对球的指向支点力记 \(N\)（拉力为正），得
\[
N=m\ell\dot\theta^2+mg\cos\theta-ma\sin\theta.
\]
初始时 \(\theta=0,\dot\theta=0\)，\(N(0)=mg\)；稳定平衡时 \(\theta=-\varphi\)、静止，\(N_s=mg_*\)。当 \(a\to0\)，\(\varphi\to0,\omega_s\to\sqrt{g/\ell},N_s\to mg\)，回到普通单摆。
